\documentclass[11pt,letterpaper]{article}
% KJO_preamble.tex
% Shared LaTeX preamble for all KJO documentation documents.
% \input this file immediately after \documentclass{}.

% --- Encoding and fonts -------------------------------------------------------
\usepackage[utf8]{inputenc}
\usepackage[T1]{fontenc}
\usepackage{lmodern}
\usepackage{microtype}

% --- Page geometry ------------------------------------------------------------
\usepackage[
  letterpaper,
  top=1in, bottom=1in,
  left=1.25in, right=1in
]{geometry}

% --- Core utilities -----------------------------------------------------------
\usepackage{amsmath}
\usepackage{graphicx}
\usepackage{xcolor}
\usepackage{parskip}        % space between paragraphs; no indent
\usepackage{enumitem}
\usepackage{caption}
\usepackage{subcaption}
\usepackage{float}

% --- Tables -------------------------------------------------------------------
\usepackage{booktabs}
\usepackage{longtable}
\usepackage{array}
\usepackage{multirow}

% --- Code listings ------------------------------------------------------------
\usepackage{listings}
\lstset{
  basicstyle=\ttfamily\small,
  keywordstyle=\bfseries\color{KJOblue},
  commentstyle=\itshape\color{gray},
  stringstyle=\color{KJOgreen!70!black},
  breaklines=true,
  frame=single,
  framesep=4pt,
  xleftmargin=6pt,
  xrightmargin=4pt,
  numbers=left,
  numberstyle=\tiny\color{gray},
  numbersep=8pt,
  language=C++,
}

% --- TikZ / PGF ---------------------------------------------------------------
\usepackage{tikz}
\usetikzlibrary{
  arrows.meta,
  shapes.geometric,
  shapes.misc,
  shapes.symbols,
  positioning,
  fit,
  calc,
  backgrounds,
  decorations.pathreplacing,
  decorations.markings,
  matrix,
}
\usepackage{pgfplots}
\pgfplotsset{compat=1.18}

% --- Headers / footers --------------------------------------------------------
\usepackage{fancyhdr}
\usepackage{lastpage}
\pagestyle{fancy}
\fancyhf{}
\renewcommand{\headrulewidth}{0.4pt}
\renewcommand{\footrulewidth}{0.4pt}
% Per-document: \lhead{}, \rhead{}, \cfoot{} are set in the main .tex file.

% --- Section formatting -------------------------------------------------------
\usepackage{titlesec}
\titleformat{\section}{\large\bfseries\color{KJOblue}}{{\thesection}}{0.8em}{}[\titlerule]
\titleformat{\subsection}{\normalsize\bfseries\color{KJOblue!80!black}}{{\thesubsection}}{0.8em}{}
\titleformat{\subsubsection}{\normalsize\bfseries}{{\thesubsubsection}}{0.8em}{}

% --- Hyperlinks ---------------------------------------------------------------
\usepackage[
  colorlinks=true,
  linkcolor=KJOblue,
  urlcolor=KJOblue!70!black,
  citecolor=KJOblue,
  pdfborder={0 0 0},
]{hyperref}

% --- Color palette -----------------------------------------------------------
\definecolor{KJOblue}   {RGB}{ 30,  90, 160}   % RCU
\definecolor{KJOred}    {RGB}{180,  30,  30}   % EMU
\definecolor{KJOgreen}  {RGB}{ 20, 120,  60}   % GSEMU
\definecolor{KJOorange} {RGB}{200, 110,  20}   % GSE / pressurization
\definecolor{KJOgray}   {RGB}{100, 100, 100}   % Buses / links
\definecolor{KJOpurple} {RGB}{110,  40, 140}   % Shared libraries
\definecolor{KJObg}     {RGB}{248, 248, 252}   % Light background for code

% --- Convenience macros -------------------------------------------------------

% Hardware component name (small caps)
\newcommand{\hw}[1]{\textsc{#1}}

% Software module name (monospace)
\newcommand{\sw}[1]{\texttt{#1}}

% Command name (monospace bold)
\newcommand{\cmd}[1]{\texttt{\textbf{#1}}}

% Pin / signal name
\newcommand{\sig}[1]{\texttt{#1}}

% Note / callout box
\usepackage{tcolorbox}
\tcbuselibrary{skins}
\newtcolorbox{notebox}{
  enhanced,
  colback=KJOblue!5,
  colframe=KJOblue!40,
  fonttitle=\bfseries,
  title=Note,
  left=4pt, right=4pt, top=2pt, bottom=2pt,
}
\newtcolorbox{warningbox}{
  enhanced,
  colback=KJOred!5,
  colframe=KJOred!40,
  fonttitle=\bfseries\color{KJOred},
  title=Warning,
  left=4pt, right=4pt, top=2pt, bottom=2pt,
}

% Title page macro: \KJOtitlepage{title}{subtitle}{version}{date}{description}
\newcommand{\KJOtitlepage}[5]{%
  \begin{titlepage}
    \centering
    \vspace*{2cm}
    {\Huge\bfseries\color{KJOblue} #1 \par}
    \vspace{0.5cm}
    {\Large\color{KJOgray} #2 \par}
    \vspace{1.5cm}
    \rule{0.5\textwidth}{0.4pt}\par
    \vspace{1cm}
    {\large Ken Overton \par}
    \vspace{0.3cm}
    {\normalsize \textit{KJO Liquid Rocket Motor Control System} \par}
    \vspace{1.5cm}
    \begin{tabular}{ll}
      \textbf{Version:} & #3 \\[4pt]
      \textbf{Date:}    & #4 \\
    \end{tabular}
    \vspace{1.5cm}\par
    \begin{minipage}{0.75\textwidth}
      \centering
      \normalsize #5
    \end{minipage}
    \vfill
    {\small\color{KJOgray} \textcopyright\ 2025--2026 Ken Overton.
     All rights reserved. \par}
  \end{titlepage}%
}

\newcommand{\docversion}{1.0}
\newcommand{\docdate}{September 2026}
\lhead{\textsc{KJO Engine Management Unit}}
\rhead{\textsc{EMU Reference} --- v\docversion}
\cfoot{\thepage\ of \pageref{LastPage}}
\hypersetup{pdftitle={KJO EMU Reference}, pdfauthor={Ken Overton}}

\begin{document}
% KJO_tikz_styles.tex
% Shared TikZ node and edge styles for all KJO block diagrams.
% \input this file inside the document body before any tikzpicture environments.

\tikzset{
  % ── Hardware block (peripheral chip / PCB component) ──────────────────────
  hwblock/.style={
    rectangle, draw=KJOgray, line width=1pt,
    fill=KJOgray!10, rounded corners=4pt,
    minimum width=2.6cm, minimum height=0.8cm,
    font=\small, align=center, inner sep=4pt,
  },
  % Sub-variant: MCU (larger, blue border)
  mcublock/.style={
    hwblock,
    draw=KJOblue, fill=KJOblue!8,
    minimum width=3.0cm, minimum height=1.0cm,
    font=\small\bfseries,
  },
  % Sub-variant: power / battery
  powerblock/.style={
    hwblock,
    draw=KJOorange, fill=KJOorange!8,
  },
  % ── Software module box ────────────────────────────────────────────────────
  swmodule/.style={
    rectangle, draw=KJOblue, line width=0.8pt,
    fill=KJOblue!6, rounded corners=3pt,
    minimum width=2.8cm, minimum height=0.7cm,
    font=\small, align=center, inner sep=3pt,
  },
  % Sub-variant: shared library
  sharedlib/.style={
    swmodule,
    draw=KJOpurple, fill=KJOpurple!6,
    dashed,
  },
  % ── Bus / link annotation ──────────────────────────────────────────────────
  busline/.style={
    draw=KJOgray, line width=1.5pt,
    -{Stealth[length=6pt,width=4pt]},
  },
  bidbusline/.style={
    draw=KJOgray, line width=1.5pt,
    {Stealth[length=6pt,width=4pt]}-{Stealth[length=6pt,width=4pt]},
  },
  radioline/.style={
    draw=KJOblue!60, line width=1.5pt, dashed,
    {Stealth[length=6pt,width=4pt]}-{Stealth[length=6pt,width=4pt]},
  },
  gpioline/.style={
    draw=KJOgreen!60!black, line width=1.2pt,
    -{Stealth[length=5pt,width=4pt]},
  },
  % ── Flowchart shapes ──────────────────────────────────────────────────────
  flowproc/.style={
    rectangle, draw=KJOblue!70, line width=0.8pt,
    fill=KJOblue!8, rounded corners=3pt,
    minimum width=3.2cm, minimum height=0.7cm,
    font=\small, align=center, inner sep=3pt,
  },
  flowdec/.style={
    diamond, draw=KJOorange, line width=0.8pt,
    fill=KJOorange!8,
    minimum width=2.2cm, minimum height=0.9cm,
    font=\small, align=center, inner sep=1pt,
  },
  flowterminal/.style={
    rounded rectangle, draw=KJOblue, line width=1pt,
    fill=KJOblue!12,
    minimum width=2.8cm, minimum height=0.7cm,
    font=\small\bfseries, align=center, inner sep=3pt,
  },
  flowarrow/.style={
    draw=KJOgray, line width=0.9pt,
    -{Stealth[length=5pt,width=4pt]},
  },
  % ── Propellant schematic styles ────────────────────────────────────────────
  % Generic pipe line
  pipe/.style={
    draw=black, line width=1.8pt,
  },
  pipearrow/.style={
    pipe,
    -{Stealth[length=7pt,width=5pt]},
  },
  % Oxidizer (N2O) — blue
  oxpipe/.style={
    draw=KJOblue, line width=1.8pt,
  },
  oxpipearrow/.style={
    oxpipe,
    -{Stealth[length=7pt,width=5pt]},
  },
  % Fuel — red
  fuelpipe/.style={
    draw=KJOred, line width=1.8pt,
  },
  fuelpipearrow/.style={
    fuelpipe,
    -{Stealth[length=7pt,width=5pt]},
  },
  % GSE fill line — green
  gsepipe/.style={
    draw=KJOgreen!70!black, line width=1.8pt,
  },
  gsepipearrow/.style={
    gsepipe,
    -{Stealth[length=7pt,width=5pt]},
  },
  % Pressurization — orange dashed
  presspipe/.style={
    draw=KJOorange, line width=1.4pt, dashed,
  },
  presspipearrow/.style={
    presspipe,
    -{Stealth[length=6pt,width=4pt]},
  },
  % Tank (tall rectangle, rounded corners)
  tank/.style={
    rectangle, draw=black, line width=1.5pt,
    rounded corners=4pt, fill=white,
    minimum width=1.8cm, minimum height=2.4cm,
    font=\small\bfseries, align=center, inner sep=4pt,
  },
  n2otank/.style={
    tank, draw=KJOblue, fill=KJOblue!8,
  },
  fueltank/.style={
    tank, draw=KJOred, fill=KJOred!8,
  },
  srctank/.style={
    tank, draw=KJOgreen!70!black, fill=KJOgreen!6,
  },
  % Valve box (servo-actuated)
  valvebox/.style={
    rectangle, draw=black, line width=1pt,
    fill=white, minimum width=2.2cm, minimum height=0.65cm,
    font=\footnotesize\bfseries, align=center, inner sep=3pt,
  },
  gseval/.style={
    valvebox, draw=KJOgreen!70!black, fill=KJOgreen!8,
  },
  emuval/.style={
    valvebox, draw=KJOred, fill=KJOred!6,
  },
  % Component (generic inline component)
  component/.style={
    rectangle, draw=KJOgray, line width=0.8pt,
    fill=white!96!gray, minimum width=2.2cm, minimum height=0.6cm,
    font=\footnotesize, align=center, inner sep=3pt,
  },
  % Instrument callout circle (PT, TT, TC)
  instrument/.style={
    circle, draw=KJOgray, line width=0.8pt,
    fill=white, minimum size=0.6cm,
    font=\tiny\bfseries, align=center, inner sep=0pt,
  },
  % Sensor lead line
  sensorlead/.style={
    draw=KJOgray, line width=0.7pt,
  },
  % Boundary box (dashed outline for GSE / Rocket regions)
  gseboundary/.style={
    draw=KJOgreen!50, line width=1pt, dashed,
    rounded corners=6pt,
  },
  rocketboundary/.style={
    draw=KJOred!40, line width=1pt,
    rounded corners=6pt,
  },
  % Sequence diagram lifeline
  lifeline/.style={
    draw=KJOgray, line width=0.6pt, dashed,
  },
  % Sequence diagram message arrow
  seqarrow/.style={
    draw=KJOblue, line width=0.9pt,
    -{Stealth[length=5pt,width=4pt]},
  },
  % State node (FSM)
  statenode/.style={
    circle, draw=KJOblue, line width=1pt,
    fill=KJOblue!10,
    minimum size=1.2cm,
    font=\small\bfseries, align=center,
  },
  statetransition/.style={
    draw=KJOgray, line width=0.9pt,
    -{Stealth[length=5pt,width=4pt]},
    font=\footnotesize,
  },
}


\KJOtitlepage%
  {Engine Management Unit (EMU)}%
  {Hardware and Software Reference}%
  {\docversion}{\docdate}%
  {This document describes the hardware, software architecture, valve control
   system, Fill valve handshake protocol, sensor subsystems, buzzer and visual
   indicator state machine, and control flow of the Engine Management Unit.
   The EMU is mounted inside the rocket and is the primary actuator and
   sensor node of the KJO system.}

\tableofcontents\newpage

%═══════════════════════════════════════════════════════════════════════════════
\section{Hardware Overview}
%═══════════════════════════════════════════════════════════════════════════════

\subsection{Component Summary}

\begin{center}
\begin{tabularx}{\textwidth}{@{}lp{6.2cm}lY@{}}
\toprule
Component & Part & Interface & Notes \\
\midrule
Microcontroller   & Adafruit Feather M0 (ATSAMD21G18A, 48~MHz)  & ---     & \\
OLED Display      & SH1107, 128$\times$64, 1.12\textquotedbl      & I\textsuperscript{2}C & \\
GPIO Expander     & MCP23X17 (16 I/O)                            & I\textsuperscript{2}C & Relays, buttons, AUX\_IO \\
PWM Servo Wing    & Adafruit 8-channel wing (PCA9685)            & I\textsuperscript{2}C & Ox, Fuel servos \\
ADC               & ADS1015 (12-bit, 4-ch)                       & I\textsuperscript{2}C & Ch0: Airstart/LCO ignition sense; Ch1: chamber pressure; Ch2: tank pressure; Ch3: battery voltage \\
Thermocouple IF   & MAX31856 (2$\times$)                         & SPI     & TC1, TC2 \\
Real-Time Clock   & PCF8523                                      & I\textsuperscript{2}C & \\
LoRa Radio        & HopeRF RFM95W (Semtech SX1276), 915~MHz      & SPI     & \\
CAN Controller    & MCP2515 (Adafruit CAN FeatherWing)           & SPI     & Link to GSEMU/LCMU \\
SD Card           & Standard SPI SD socket                       & SPI     & \\
NeoPixel          & WS2812B 4$\times$8 array                     & GPIO    & Status indicator \\
Igniter Relay     & SPDT relay via MCP23X17                      & GPIO    & \\
Buzzer            & Passive piezo via MCP23X17 relay             & GPIO    & \\
Battery           & 3.7~V LiPo; voltage divider on ADS1015       & ADC     & \\
\bottomrule
\end{tabularx}
\end{center}

\subsection{MCP23X17 GPIO Expander Pin Map}

\begin{center}
\begin{tabularx}{\textwidth}{@{}llY@{}}
\toprule
Pin Name & Direction & Function \\
\midrule
\sig{BUTTON\_A\_EIO}                              & IN  & Front-panel button A (hold-to-open Ox valve, dumps via chamber) \\
\sig{BUTTON\_B\_EIO}                              & IN  & Front-panel button B; currently unused (see §\ref{sec:emu-local-buttons}) \\
\sig{BUTTON\_C\_EIO}                              & IN  & Front-panel button C; currently unused (see §\ref{sec:emu-local-buttons}) \\
\sig{IGNITER\_EIO}                                & OUT & Igniter relay drive \\
\sig{BUZZER\_EIO}                                 & OUT & Buzzer relay drive \\
\sig{AUX\_IO\_1} (\sig{EMU\_LINK\_SENSE\_EIO}) & IN (INPUT\_PULLUP) & Umbilical link-disconnect sense; wired to GSEMU chassis GND (LOW = connected, HIGH = separated). Fill valve CMD/STATUS formerly carried here has moved to CAN. \\
\sig{AUX\_IO\_2}                                & --- & Retired/unused (formerly \sig{FILL\_VALVE\_STATUS\_EIO}). \\
\sig{AUX\_IO\_3} (\sig{QR\_SENSE\_EIO})         & IN (INPUT\_PULLUP) & Umbilical QR physical-connection sense; wired to GSEMU chassis GND (LOW = connected, HIGH = separated). QR release itself is commanded over CAN (\cmd{CAN\_QR\_RELEASE}), not this line. \\
\sig{AUX\_IO\_4}                                & --- & Retired/unused (formerly \sig{RELEASE\_STATE\_EIO}; its sense function consolidated onto \sig{QR\_SENSE\_EIO}). \\
\sig{AUX\_IO\_5}                                & --- & Retired/unused (formerly \sig{LAUNCH\_ENABLE\_EIO}). \\
\sig{AUX\_IO\_6}                                & --- & Retired/unused (formerly \sig{REMOTE\_START\_EIO}). \\
\bottomrule
\end{tabularx}
\end{center}

\subsection{Hardware Block Diagram}

\begin{figure}[H]
\centering
\begin{tikzpicture}[node distance=0.9cm and 1.8cm]
  \node[mcublock, minimum width=3.0cm, minimum height=1.1cm] (mcu)
    {Feather M0\\(ATSAMD21G18A)};
  % I2C left
  \node[hwblock, left=2.3cm of mcu, yshift= 3.2cm, minimum width=3.2cm] (oled) {SH1107 OLED};
  \node[hwblock, left=2.3cm of mcu, yshift= 2.0cm, minimum width=3.2cm] (gpio) {MCP23X17\\GPIO Exp.};
  \node[hwblock, left=2.3cm of mcu, yshift= 0.8cm, minimum width=3.2cm] (pwm)  {PCA9685\\Servo Wing};
  \node[hwblock, left=2.3cm of mcu, yshift=-0.4cm, minimum width=3.2cm] (adc)  {ADS1015 ADC};
  \node[hwblock, left=2.3cm of mcu, yshift=-1.6cm, minimum width=3.2cm] (rtc)  {PCF8523 RTC};
  % SPI right
  \node[hwblock, right=2.3cm of mcu, yshift= 2.0cm, minimum width=3.0cm] (rf)  {RFM95W LoRa};
  \node[hwblock, right=2.3cm of mcu, yshift= 0.8cm, minimum width=3.0cm] (tc1) {MAX31856 TC×2};
  \node[hwblock, right=2.3cm of mcu, yshift=-0.4cm, minimum width=3.0cm] (sd)  {SD Card};
  \node[hwblock, right=2.3cm of mcu, yshift=-1.6cm, minimum width=3.0cm] (can) {MCP2515\\CAN FeatherWing};
  % GPIO below
  \node[hwblock, below=1.5cm of mcu, xshift=-1.8cm, minimum width=2.8cm] (neo) {NeoPixel 4×8};
  \node[hwblock, below=1.5cm of mcu, xshift= 1.8cm, minimum width=2.8cm] (buz) {Buzzer};
  % Power
  \node[powerblock, below=3.0cm of mcu, minimum width=3.0cm] (bat) {LiPo Battery};
  % I2C bus (left)
  \foreach \nd in {oled,gpio,pwm,adc,rtc}
    \draw[bidbusline] (mcu.west |- \nd) -- (\nd.east)
      node[midway,above,font=\tiny]{I\textsuperscript{2}C};
  % SPI bus (right)
  \foreach \nd in {rf,tc1,sd,can}
    \draw[bidbusline] (mcu.east |- \nd) -- (\nd.west)
      node[midway,above,font=\tiny]{SPI};
  % GPIO
  \draw[busline] (mcu.south -| neo) -- (neo.north);
  \draw[busline] (mcu.south -| buz) -- (buz.north);
  % Power
  \draw[busline,draw=KJOorange] (bat.north) -- ++(0,0.6) -| (mcu.south)
    node[near end,right,font=\tiny]{3.7V};
\end{tikzpicture}
\caption{EMU hardware block diagram.}
\end{figure}

%═══════════════════════════════════════════════════════════════════════════════
\section{Software Architecture}
%═══════════════════════════════════════════════════════════════════════════════

\subsection{Module Summary}

\begin{center}
\begin{tabularx}{\textwidth}{@{}llY@{}}
\toprule
Module & Location & Responsibility \\
\midrule
\sw{main.cpp}              & local  & Setup, loop, command dispatch, state machines \\
\sw{KJO\_Valve}            & local  & Base valve: servo drive, pot feedback \\
\sw{KJO\_QR\_Master}       & local  & Umbilical connection-sense reader (see §6) \\
\sw{KJO\_Control\_System}  & local  & \sw{Execute\_Command()} routing table \\
\sw{KJO\_Pressure}         & local  & Pressure transducer read / convert \\
\sw{KJO\_Thermocouple}     & local  & MAX31856 thermocouple read (TC1/TC2 -- all temperature sensing on this board) \\
\sw{KJO\_Analog}           & local  & ADS1015 ADC wrapper, Airstart/LCO channel \\
\sw{KJO\_GPIO}             & local  & MCP23X17 pin definitions, Relay class \\
\sw{KJO\_EMU\_Display}     & local  & SD log-file finder (\sw{Find\_Available\_File()}) only -- OLED display moved to shared \sw{KJO\_Status\_Display} \\
\sw{KJO\_Indicator}        & local  & NeoPixel status indicator (solid, timed, flash) \\
\sw{KJO\_Logging}          & local  & SD log file constants \\
\sw{KJO\_Units}            & local  & Unit conversion helpers (\sw{C\_To\_F()}) \\
\sw{KJO\_Burn\_Buffer}     & local (\texttt{lib/}) & Pre/post-trigger RAM ring buffer for burn data capture \\
\sw{KJO\_Radio}            & shared & Radio pin constants, frequency \\
\sw{KJO\_Command\_Defs}    & shared & LoRa command codes, packet structures \\
\sw{KJO\_CAN\_Command\_Defs} & shared & CAN command codes, packet structures, node IDs \\
\sw{KJO\_CAN\_Client}      & shared & CAN transport (send/poll/request-with-timeout) \\
\sw{KJO\_System\_Config}   & shared & Version, platform constants \\
\sw{KJO\_Status} / \sw{KJO\_Status\_Display} & shared & Tagged Serial/SD logging (\sw{Log\_Message()}) and scrolling OLED status (\sw{Report\_Status()}) \\
\bottomrule
\end{tabularx}
\end{center}

\subsection{Main Loop Control Flow}

\begin{figure}[H]
\centering
\resizebox{!}{0.92\textheight}{%
\begin{tikzpicture}[node distance=0.28cm]
  \node[flowterminal] (start)  {loop()};
  \node[flowproc, below=of start]  (rel)  {Service relay timers\\(Igniter\_Relay, Buzzer)};
  \node[flowproc, below=of rel]    (ind)  {Indicator.update()};
  \node[flowproc, below=of ind]    (btn)  {Check\_Buttons()};
  \node[flowproc, below=of btn]    (vchk) {Check\_Valve\_Move()};
  \node[flowproc, below=of vchk]   (bcan) {Check\_Boot\_CAN\_Sequence()};
  \node[flowproc, below=of bcan]   (pping)  {Check\_Presence\_Ping()};
  \node[flowproc, below=of pping]  (fpoll)  {Check\_Fill\_Poll()};
  \node[flowproc, below=of fpoll]  (eng)  {Check\_Engine\_Start()};
  \node[flowproc, below=of eng]    (abrt) {Check\_Abort\_Sequence()};
  \node[flowproc, below=of abrt]   (air)  {Check\_Airstart\_Input()};
  \node[flowproc, below=of air]    (qrcs) {Check\_QR\_Connection\_Status()};
  \node[flowproc, below=of qrcs]   (pwarn) {Check\_Pressure\_Warning()};
  \node[flowproc, below=of pwarn]  (oxover) {Check\_Ox\_Tank\_Overpressure()};
  \node[flowproc, below=of oxover] (prec) {Check\_Pressure\_Recording()};
  \node[flowproc, below=of prec]   (burn) {Check\_Burn\_Recording()};
  \node[flowproc, below=of burn]   (los)  {Check\_Link\_LOS()};
  \node[flowproc, below=of los]    (iab)  {Check\_Indicator\_And\_Buzzer()};
  \node[flowproc, below=of iab]    (rx)   {Check\_Comms()};
  \node[flowproc, below=of rx]     (can)  {Check\_CAN()};
  \draw[flowarrow] (start)--(rel);
  \draw[flowarrow] (rel)--(ind);
  \draw[flowarrow] (ind)--(btn);
  \draw[flowarrow] (btn)--(vchk);
  \draw[flowarrow] (vchk)--(bcan);
  \draw[flowarrow] (bcan)--(pping);
  \draw[flowarrow] (pping)--(fpoll);
  \draw[flowarrow] (fpoll)--(eng);
  \draw[flowarrow] (eng)--(abrt);
  \draw[flowarrow] (abrt)--(air);
  \draw[flowarrow] (air)--(qrcs);
  \draw[flowarrow] (qrcs)--(pwarn);
  \draw[flowarrow] (pwarn)--(oxover);
  \draw[flowarrow] (oxover)--(prec);
  \draw[flowarrow] (prec)--(burn);
  \draw[flowarrow] (burn)--(los);
  \draw[flowarrow] (los)--(iab);
  \draw[flowarrow] (iab)--(rx);
  \draw[flowarrow] (rx)--(can);
  \draw[flowarrow] (can.south) -- ++(0,-0.4) -- ++(4.5,0)
    -- ++(0,22.0) -- (start.east);
\end{tikzpicture}%
}
\caption{EMU main loop control flow. \sw{ADC\_Diagnostic\_Print()} (bench-only, compiled out by default) runs after \sw{Check\_CAN()} when \texttt{ADC\_DIAGNOSTIC\_MODE} is defined -- omitted above.}
\end{figure}

%═══════════════════════════════════════════════════════════════════════════════
\section{Valve Control Subsystem}
%═══════════════════════════════════════════════════════════════════════════════

\subsection{Valve Base Class (\sw{KJO\_Valve})}

The \sw{Valve} class encapsulates a servo-driven, potentiometer-feedback valve:

\begin{itemize}[nosep]
  \item \textbf{Drive}: PWM signal to PCA9685 servo wing channel.
  \item \textbf{Feedback}: Potentiometer on ADS1015 ADC channel; position
        expressed as 0--100\% open.
  \item \textbf{Position control}: \sw{openTo(\%)} sets target; \sw{update()}
        drives servo until pot reading reaches target within dead-band.
  \item \textbf{State queries}: \sw{isMoving()}, \sw{isStopped()},
        \sw{isOpen()}, \sw{isClosed()}.
\end{itemize}

\subsection{EMU State Machine}

The EMU operates as a top-level state machine with five states:

\begin{figure}[H]
\centering
\begin{tikzpicture}[node distance=3.0cm]
  \node[statenode] (idle)     {VALVE\\IDLE};
  \node[statenode, right=of idle]  (moving)   {VALVE\\MOVING};
  \node[statenode, below=2.5cm of idle] (starting) {ENGINE\\STARTING};
  \node[statenode, right=of starting] (aborting) {ABORTING};
  \node[statenode, below=2.5cm of aborting] (latched) {ABORT\\LATCHED};

  % IDLE → MOVING
  \draw[statetransition, bend left=15]
    (idle) edge node[above, font=\tiny, align=center]{valve cmd\\received}
    (moving);
  % MOVING → IDLE
  \draw[statetransition, bend left=15]
    (moving) edge node[below, font=\tiny, align=center]{move complete\\or timeout}
    (idle);
  % MOVING self-loop
  \draw[statetransition, loop above, looseness=5]
    (moving) edge node[above, font=\tiny]{still moving}
    (moving);
  % IDLE → ENGINE_STARTING
  \draw[statetransition, bend right=10]
    (idle) edge node[left, font=\tiny, align=center]{START\_ENGINE\\or LCO/Airstart}
    (starting);
  % ENGINE_STARTING → IDLE
  \draw[statetransition, bend right=10]
    (starting) edge node[right, font=\tiny, align=center]{seq complete}
    (idle);
  % any state -> ABORTING
  \draw[statetransition]
    (starting) edge node[above, font=\tiny, align=center]{RESET (any state)}
    (aborting);
  % ABORTING -> ABORT_LATCHED
  \draw[statetransition]
    (aborting) edge node[right, font=\tiny, align=center]{tank vented\\below threshold}
    (latched);
\end{tikzpicture}
\caption{EMU top-level state machine.  HEARTBEAT is processed in any
         state; RESET is processed in any state except \texttt{ABORT\_LATCHED}
         and transitions to \texttt{ABORTING} (see §3.4 Abort Sequence).
         All other commands are dropped while \texttt{VALVE\_MOVING},
         \texttt{ENGINE\_STARTING}, \texttt{ABORTING}, or
         \texttt{ABORT\_LATCHED}. \texttt{ABORT\_LATCHED} is terminal --
         there is no software path back to \texttt{VALVE\_IDLE}.}
\end{figure}

\subsection{Engine Start Sequence (10 Steps)}

When \cmd{START\_ENGINE} or a LCO panel button press is received, the EMU
enters \texttt{ENGINE\_STARTING} and advances through ten steps (0--9)
non-blockingly in \sw{Check\_Engine\_Start()}. Propellant flow is
established (steps 6--8) \emph{before} the igniter fires (step 9) --
reversed from an earlier design that fired the igniter first. Fuel opens
before Ox since it has the longer physical path to the chamber and needs
the head start.

\begin{center}
\begin{tabular}{@{}clp{8.5cm}@{}}
\toprule
Step & Trigger & Action \\
\midrule
0 & immediate       & Verify Ox and Fuel valves are closed; abort ($r\_val=-1$) on failure \\
1 & step 0 pass     & If GSEMU is CAN-reachable, close the Fill valve; otherwise skip to step 3 \\
2 & timer / CAN ack & Wait for the Fill valve close to confirm; abort ($r\_val=-4$) on failure \\
3 & step 1/2 pass   & Send \cmd{CAN\_QR\_RELEASE} (fire-and-forget); start \texttt{QR\_RELEASE\_TIMEOUT\_MS} separation timer \\
4 & timer           & Poll \sw{QR\_Release.isSeparated()}; abort ($r\_val=-2$) on timeout \\
5 & separation      & Fixed \texttt{START\_DWELL\_MS} pause; then open Fuel feed valve \\
6 & Fuel moving     & Wait for Fuel valve to confirm open; abort via \sw{Execute\_Abort()} ($r\_val=-5$) if stuck/stalled; then start \texttt{FUEL\_FLOW\_DELAY\_MS} timer \\
7 & delay           & Open Ox feed valve \\
8 & Ox moving       & Wait for Ox valve to confirm open; abort via \sw{Execute\_Abort()} ($r\_val=-5$) if stuck/stalled; then start \texttt{OX\_FLOW\_DELAY\_MS} timer \\
9 & delay           & Fire igniter relay (\texttt{IGNITER\_FIRE\_DURATION\_MS} pulse); send OK response; set \sw{engine\_running = true} \\
\bottomrule
\end{tabular}
\end{center}

Any abort at steps 0--4 sends a \texttt{Command\_Return\_Packet} with
\texttt{status = false} and an \texttt{r\_val} error code
($-1$ = precondition failure, $-4$ = Fill valve close failure, $-2$ = QR
separation timeout). A stuck/stalled Fuel or Ox valve at steps 6 or 8
($r\_val=-5$) routes through \sw{Execute\_Abort()} instead, since
propellant may already be flowing by that point.
\cmd{RESET} cancels the sequence at any step.

\subsection{Abort Sequence}

\cmd{RESET} (or a link-loss abort in \texttt{EMU\_TEST\_MODE} -- see
\sw{Check\_Link\_LOS()}) calls \sw{Execute\_Abort()}, which performs
several actions immediately -- de-energizes the igniter, stops any
in-progress start/valve-move sequence, clears ignition-source arming,
silences the buzzer alarm -- \textbf{but deliberately leaves sensor
recording running} (it stops on its own once pressure settles). It then
arms a 7-step sequenced valve/vent procedure, advanced non-blockingly by
\sw{Check\_Abort\_Sequence()} every \sw{loop()} while \texttt{emu\_state}
is \texttt{ABORTING}:

\begin{center}
\begin{tabular}{@{}clp{8.5cm}@{}}
\toprule
Step & Trigger & Action \\
\midrule
0 & immediate     & If GSEMU is CAN-reachable, close the Fill valve; otherwise skip to step 2 \\
1 & timer / CAN ack & Wait for the Fill valve close to confirm; proceed regardless of outcome \\
2 & step 0/1 pass & Close the Fuel valve; start a confirm timeout \\
3 & timer or stopped & Wait for the Fuel valve to stop moving; proceed regardless of outcome \\
4 & step 3 pass   & Open the Ox valve (dump); start a confirm timeout \\
5 & timer or stopped & Wait for the Ox valve to stop moving; proceed regardless of outcome \\
6 & step 5 pass   & Wait (no timeout) for tank gauge pressure to drop below \texttt{ABORT\_VENTED\_THRESHOLD\_PSI}; once vented, latch \texttt{emu\_state = ABORT\_LATCHED} \\
\bottomrule
\end{tabular}
\end{center}

Steps 0--5 each proceed on their own timeout even if a confirmation never
arrives -- venting the flight-side propellants must never stall on a lost
CAN ack or a valve that fails to confirm. Step 6 has no timeout: a stuck
pressure sensor or a genuine slow leak both mean the system is not yet
safe, and there is no further fallback to proceed to.
\texttt{ABORT\_LATCHED} is terminal -- \sw{Execute\_Abort()} itself
refuses re-entry once latched, and there is no software path back to
\texttt{VALVE\_IDLE}.

A separate, unconditional safety poll (\sw{Check\_Ox\_Tank\_Overpressure()},
every \sw{loop()} regardless of state) calls \sw{Execute\_Abort()}
automatically if Ox tank pressure stays at or above
\texttt{MAX\_OX\_TANK\_PRESSURE} for \texttt{OX\_TANK\_PRESSURE\_TIME\_MS}
(see §10 Parameters \& Flags) -- an abort does not require an explicit
\cmd{RESET} command.

\subsection{Servo Mechanical Geometry and Calibration}

\subsubsection*{Gear Train and Servo Travel}

Each valve is actuated by a ServoCity 2000-0025-0504 5-turn servo through a
16-tooth pinion driving a 60-tooth valve gear:

\begin{center}
\begin{tabular}{@{}ll@{}}
\toprule
Parameter & Value \\
\midrule
Gear ratio                             & $60 \div 16 = \mathbf{3.75:1}$ \\
Valve travel (open $\to$ close)        & 90\textdegree \\
Required servo rotation                & $90 \times 3.75 = 337.5$\textdegree \\
Servo full mechanical range            & 1800\textdegree\ (5 turns) over 500--2500~$\mu$s \\
PCA9685 frequency                      & 1100~Hz (non-standard; improves count resolution) \\
PWM period at 1100~Hz                  & $1 \div 1100 = 909.09~\mu$s \\
$\mu$s per 12-bit count                & $909.09 \div 4096 \approx 0.2219~\mu$s/count \\
$\mu$s for 337.5\textdegree\ servo travel& $337.5 \times (2000~\mu\text{s} \div 1800\textdegree) = 375~\mu$s \\
\textbf{Theoretical $\Delta$ counts}   & $375 \div 0.2219 \approx \mathbf{1690}$ counts \\
\bottomrule
\end{tabular}
\end{center}

\begin{notebox}
Running the PCA9685 at 1100~Hz (instead of the standard 50~Hz) compresses the
usable pulse-width span into the lower portion of the 12-bit count range,
effectively improving resolution over the $\sim$1.1 turns of servo travel
actually used.
\end{notebox}

\subsubsection*{Empirical Calibration (March 2026)}

Bench testing showed the theoretical 1690-count span caused both of EMU's
valves (Ox and Fuel -- the Fill valve is GSEMU's actuator, not EMU's; see
its own reference doc for its calibration) to travel noticeably \emph{more}
than 90\textdegree.  Both were re-calibrated.  The empirical open-to-close
deltas average $\approx$1645~counts ($\sim$3\% below theoretical),
consistent with the servo being slightly more sensitive in the operating
band than the linear model predicts.

\begin{center}
\begin{tabular}{@{}lrrrrrrrr@{}}
\toprule
Valve & \multicolumn{4}{c}{Servo PWM (12-bit counts)} & \multicolumn{3}{c}{Pot ADC (12-bit counts)} \\
\cmidrule(lr){2-5}\cmidrule(lr){6-8}
 & Open & Close & $\Delta$ & Implied travel & Open & Closed & $\Delta$ \\
\midrule
Oxidizer & 2050 & 3690 & 1640 & $\sim$87.4\textdegree & 1139 & 2857 & 1718 \\
Fuel     & 2130 & 3780 & 1650 & $\sim$87.9\textdegree & 1170 & 2918 & 1748 \\
\bottomrule
\end{tabular}
\end{center}

\subsection{Local Button Functions}
\label{sec:emu-local-buttons}

\sw{Check\_Buttons()} polls the three front-panel buttons (MCP23X17
expander) every \sw{loop()} iteration, but is skipped entirely while
\texttt{emu\_state} is \texttt{ENGINE\_STARTING}, \texttt{VALVE\_MOVING},
\texttt{ABORTING}, or \texttt{ABORT\_LATCHED}.

\begin{center}
\begin{tabularx}{\textwidth}{@{}lY@{}}
\toprule
Button & Action \\
\midrule
A & \textbf{Hold-to-open}: opens the Ox valve (dumps pressurant through the
    chamber) while held; closes it on release. One short beep on press. \\
B & Currently unused. Formerly commanded the Fill valve closed directly;
    the Fill valve is now GSEMU's own actuator, commanded only over CAN
    (\cmd{BEGIN\_FILL}/\cmd{STOP\_FILL}). Local bench testing of the Fill
    valve is done via GSEMU's own front-panel buttons. \\
C & Currently unused, for the same reason as Button~B. \\
\bottomrule
\end{tabularx}
\end{center}

\begin{warningbox}
Button~A opens the Ox valve immediately while held, venting
N\textsubscript{2}O through the combustion chamber. Ensure the area is
clear before pressing it.
\end{warningbox}

%═══════════════════════════════════════════════════════════════════════════════
\section{Fill Valve CAN Protocol}
%═══════════════════════════════════════════════════════════════════════════════

The Fill valve servo and potentiometer reside on the GSEMU. The wired
\sig{AUX\_IO\_1}/\sig{AUX\_IO\_2} CMD/STATUS handshake this section
formerly described has been retired; the EMU now controls the Fill
valve entirely over CAN.

\textbf{Open/close sequence:}
\begin{sloppypar}
\begin{enumerate}[nosep]
  \item EMU sends \cmd{CAN\_OPEN\_FILL\_VALVE} or \cmd{CAN\_CLOSE\_FILL\_VALVE}
        to \texttt{CAN\_NODE\_GSEMU} and awaits a response.
  \item GSEMU's \sw{Check\_CAN()} calls \sw{Fill\_Valve.open()}/\sw{close()}
        immediately, but \emph{defers} its CAN response.
  \item GSEMU's \sw{Check\_Fill\_Valve\_Move()} (called every \sw{loop()})
        sends the deferred response once \sw{Fill\_Valve.isStopped()}
        confirms the move (or its own \texttt{MOVE\_TIMEOUT}-based deadline
        expires), reporting the resulting position.
  \item EMU treats the request as failed if no response arrives within
        \texttt{EMU\_GSEMU\_FILL\_VALVE\_MOVE\_TIMEOUT\_MS} (see §10
        Parameters \& Flags), which is sized to cover GSEMU's own
        move watchdog plus CAN round-trip margin.
\end{enumerate}
\end{sloppypar}

%═══════════════════════════════════════════════════════════════════════════════
\section{Command Dispatch}
%═══════════════════════════════════════════════════════════════════════════════

\sw{Check\_Comms()} receives packets from the RCU and routes them.
Commands are processed in priority order:

\begin{sloppypar}
\begin{enumerate}[nosep]
  \item \cmd{HEARTBEAT} — always processed; returns
        \texttt{Heartbeat\_Response\_Packet} with EMU voltage and
        the Ox/Fuel valve positions (\texttt{fill\_pct} is hardcoded 0 --
        Fill state is polled via CAN/\cmd{READ\_VALVE\_POSITION}, not
        heartbeat).
  \item \cmd{RESET} — always processed except in \texttt{ABORT\_LATCHED};
        calls \sw{Execute\_Abort()} (see §3.4 Abort Sequence).
  \item All other commands are \emph{dropped} if \texttt{emu\_state} is
        \texttt{VALVE\_MOVING}, \texttt{ENGINE\_STARTING},
        \texttt{ABORTING}, or \texttt{ABORT\_LATCHED}.
  \item Per-command interceptors (deferred valve moves, start sequence,
        recording, CAN relays to GSEMU/LCMU).
  \item All remaining commands dispatched to \sw{Execute\_Command()}.
\end{enumerate}
\end{sloppypar}

\subsection{Command Table}

\begin{sloppypar}
Codes 1--35 are production commands; codes 36--39 exist only for
RCU\_UNIT\_TEST, a separate bench-test harness firmware with an expanded
command set for exercising individual actuators directly -- production
RCU never sends them.
\end{sloppypar}

\begingroup
\footnotesize
\begin{longtable}{@{}rlp{8.2cm}@{}}
\toprule
Code & Command & EMU Action \\
\midrule
1 & \cmd{BEGIN\_FILL} & Open Fill valve (CAN to GSEMU); start continuous pressure recording \\
2 & \cmd{STOP\_FILL} & Close Fill valve; recording continues until ambient \\
3 & \cmd{OPEN\_OXIDIZER\_FEED} & Open Ox valve (deferred) \\
4 & \cmd{CLOSE\_OXIDIZER\_FEED} & Close Ox valve (deferred) \\
5 & \cmd{OPEN\_FUEL\_FEED} & Open Fuel valve (deferred) \\
6 & \cmd{CLOSE\_FUEL\_FEED} & Close Fuel valve (deferred) \\
7 & \cmd{OXIDIZER\_FLOW\_PROP} & Dead: case exists in \sw{Execute\_Command()} but does nothing \\
8 & \cmd{FUEL\_FLOW\_PROP} & Dead: case exists in \sw{Execute\_Command()} but does nothing \\
9 & \cmd{DUMP} & Open Ox valve to vent pressurant through the chamber (deferred; sends response when complete; no dedicated Dump valve) \\
10 & \cmd{START\_ENGINE\_FLOWS} & Close Fill; open Ox + Fuel simultaneously (deferred dual-valve) \\
11 & \cmd{STOP\_ENGINE\_FLOWS} & Close Ox + Fuel simultaneously (deferred dual-valve) \\
12 & \cmd{ARM\_LAUNCH} & Arms/disarms the currently-selected ignition source (\texttt{param} $\neq$ 0 = arm). Fails if arming with no source selected. If source is LCO, relays \cmd{CAN\_ARM\_LCO\_WATCH} to GSEMU (response deferred); RCU/Airstart sources arm/disarm immediately, purely locally. Renamed from \cmd{ENABLE\_LCO\_CONTROL} -- see §10 Version History \\
13 & \cmd{START\_ENGINE} & Fails immediately unless armed for the RCU source; otherwise arms the 10-step ignition sequence (response sent asynchronously when it completes) \\
14 & \cmd{RESET} & \sw{Execute\_Abort()}: igniter off, ignition disarmed, buzzer alarm silenced, recording left running; arms the 7-step sequenced Fill/Fuel-close, Ox-open, vent-and-latch procedure (see §3.4) \\
15 & \cmd{READ\_TANK\_PRESSURE} & Return tank pressure (PSI) \\
16 & \cmd{READ\_CHAMBER\_PRESSURE} & Return chamber pressure (PSI) \\
17 & \cmd{READ\_TEMP2} & Return TC2 thermocouple temperature (formerly named \cmd{READ\_TANK\_TEMP}) \\
18 & \cmd{READ\_TEMP1} & Return TC1 thermocouple temperature (formerly named \cmd{READ\_CHAMBER\_TEMP}) \\
19 & \cmd{START\_PRESSURE\_RECORDING} & Activate continuous sensor recording \\
20 & \cmd{STOP\_PRESSURE\_RECORDING} & Deactivate continuous sensor recording \\
21 & \cmd{READ\_VALVE\_POSITION} & Return current position (\%) of the valve specified in \texttt{param} \\
22 & \cmd{HEARTBEAT} & Respond with EMU voltage, Ox/Fuel positions, QR/GSEMU/LCMU reachability (\texttt{fill\_pct} always 0 -- see above) \\
23 & \cmd{RCU\_STATUS} & Relay: query GSEMU health (\cmd{CAN\_GET\_HEALTH\_STATUS}) then LCMU health, always both, no short-circuit; response relayed back over LoRa \\
24 & \cmd{STOP\_DUMP} & Close Ox valve (deferred) \\
25 & \cmd{QR\_RELEASE} & Bench-test only: reversible timed release via \cmd{CAN\_QR\_TEST\_RELEASE} (GSEMU's local-override mechanism); production launch release happens automatically as Step 3 of the Engine Start sequence via \cmd{CAN\_QR\_RELEASE} \\
26 & \cmd{QUERY\_FILL\_PROGRESS} & Answered locally from \sw{Check\_Fill\_Poll()}'s cached state (no CAN round-trip): fill state, current weight, clamped percent, and an estimated seconds-remaining (once enough time/weight has accrued to estimate a rate) \\
27 & \cmd{TARE\_LCMU} & Relay to LCMU: \cmd{CAN\_TARE} \\
28 & \cmd{GET\_LCMU\_BATTERY} & Relay to LCMU: \cmd{CAN\_GET\_BATTERY\_VOLTAGE} \\
29 & \cmd{REPORT\_LCMU\_WEIGHT} & Relay to LCMU: \cmd{CAN\_REPORT\_CURRENT\_WEIGHT} \\
30 & \cmd{BEGIN\_WEIGHT\_RECORDING} & Relay to LCMU: \cmd{CAN\_BEGIN\_WEIGHT\_RECORDING} \\
31 & \cmd{STOP\_WEIGHT\_RECORDING} & Relay to LCMU: \cmd{CAN\_STOP\_WEIGHT\_RECORDING} \\
32 & \cmd{BEGIN\_THRUST\_RECORDING} & Relay to LCMU: \cmd{CAN\_BEGIN\_THRUST\_RECORDING} \\
33 & \cmd{STOP\_THRUST\_RECORDING} & Relay to LCMU: \cmd{CAN\_STOP\_THRUST\_RECORDING} \\
34 & \cmd{SET\_IGNITION\_SOURCE} & Selects the active ignition source (\texttt{param}: 0=RCU, 1=LCO, 2=Airstart); rejected while armed \\
35 & \cmd{ENGINE\_START\_ARMED} & EMU$\to$RCU, unprompted -- sent the instant any source arms the 10-step start sequence; \texttt{r\_val} = ignition source (same encoding as \cmd{SET\_IGNITION\_SOURCE}) \\
36 & \cmd{FIRE\_IGNITER} & RCU\_UNIT\_TEST only: direct \sw{Igniter\_Relay.on()}, no sequence/interlock \\
37 & \cmd{OPEN\_FILL\_VALVE} & RCU\_UNIT\_TEST only: direct Fill valve open, relayed as \cmd{CAN\_OPEN\_FILL\_VALVE}, bypassing the fill-target state machine entirely \\
38 & \cmd{CLOSE\_FILL\_VALVE} & RCU\_UNIT\_TEST only: direct Fill valve close, relayed as \cmd{CAN\_CLOSE\_FILL\_VALVE} \\
39 & \cmd{EXTENDED\_HEARTBEAT} & RCU\_UNIT\_TEST only: gathers GSEMU/LCMU sensor values over a multi-step CAN chain and returns \texttt{Extended\_Heartbeat\_Response\_Packet} instead of the normal heartbeat response \\
\bottomrule
\end{longtable}
\endgroup

%═══════════════════════════════════════════════════════════════════════════════
\section{Umbilical Quick-Release (\sw{QR\_Master})}
%═══════════════════════════════════════════════════════════════════════════════

\subsection{Overview}

The umbilical QR servo and its drive latch reside on the GSEMU.
EMU-side \sw{QR\_Master} no longer commands it via a dedicated wire --
release is commanded over CAN (\cmd{CAN\_QR\_RELEASE}, fire-and-forget).
\sw{QR\_Master}'s only remaining job is reading the umbilical's physical
connection-sense line:

\begin{itemize}[nosep]
  \item \sig{AUX\_IO\_3} (\sig{QR\_SENSE\_EIO}) --- INPUT\_PULLUP.  Wired
        to GSEMU chassis GND through the AUX connector.  Reads LOW while the
        umbilical is physically connected; floats HIGH when the connector is
        removed (separation confirmed).  Kept as a dedicated wire (not CAN)
        because once the umbilical separates, EMU and GSEMU can no longer
        communicate over CAN at all.
\end{itemize}

\begin{sloppypar}
\sw{QR\_Master} exposes \sw{isConnected()} and \sw{isSeparated()} reads of
this line -- it has no \sw{release()}/\sw{hold()} methods.  Release is a
one-shot CAN command sent directly from \sw{main.cpp}:

\begin{itemize}[nosep]
  \item \textbf{Production launch:} \cmd{CAN\_QR\_RELEASE} is sent
        automatically as Step 3 of the Engine Start sequence (see §3.3).
        GSEMU latches this permanently -- once triggered, the spring-loaded
        latch cannot be electrically re-engaged; only physical re-latching
        (or GSEMU's own front-panel Button~A override) changes it.
  \item \textbf{Bench testing:} RCU\_UNIT\_TEST's \cmd{QR\_RELEASE} command
        instead sends \cmd{CAN\_QR\_TEST\_RELEASE}, which drives GSEMU's
        reversible local-override mechanism (auto-reverting after a fixed
        duration) and never touches the permanent production latch.
\end{itemize}
\end{sloppypar}

\begin{warningbox}
Ensure the launch sequence has been verified and all personnel are clear
before the Engine Start sequence reaches Step 3.  The umbilical
disconnection releases pressurised N\textsubscript{2}O, and production
release cannot be electrically undone once commanded.
\end{warningbox}

%═══════════════════════════════════════════════════════════════════════════════
\section{Buzzer and Visual Indicator}
%═══════════════════════════════════════════════════════════════════════════════

\subsection{Startup Self-Test}

At power-on, the buzzer fires \textbf{five short beeps} (75~ms on / 75~ms gap)
to confirm the buzzer and relay circuits are functional.  The NeoPixel cycles
red~$\to$ yellow~$\to$ green~$\to$ off as a pixel self-test.

\subsection{NeoPixel Indicator State Machine}

\sw{Check\_Indicator\_And\_Buzzer()} runs every loop iteration and updates the
32-pixel NeoPixel array to reflect the highest-priority active condition.
Ten internal states are evaluated in descending priority (only the
indicator changes when the desired state changes from the previous tick):

\begin{center}
\begin{tabularx}{\textwidth}{@{}lYp{2.7cm}l@{}}
\toprule
Priority & Condition & Pattern & Colour \\
\midrule
1 & \texttt{emu\_state} = \texttt{ABORT\_LATCHED} & Alternate periodically$^{*}$ & Red / Yellow \\
2 & \texttt{emu\_state} = \texttt{ABORTING} & Solid & Red \\
3 & \sw{engine\_running} = true & Off (all pixels dark) & --- \\
4 & \texttt{emu\_state} = \texttt{ENGINE\_STARTING}, pre-flow phase (steps 0--5), RCU/LCO source & Solid & Green \\
5 & \texttt{emu\_state} = \texttt{ENGINE\_STARTING}, otherwise (Airstart source, or flow already started) & Off & --- \\
6 & Fill session active and in progress & Solid & Blue \\
7 & Fill session active and paused & Flash 250/750~ms & Blue \\
8 & \sw{tank\_pressurized} = true & Flash 300/500~ms & Red \\
9 & \sw{link\_connected} = true & Flash 225/4775~ms & Green \\
10 & (unlinked) & Flash 225/4775~ms & Yellow \\
\bottomrule
\end{tabularx}
\end{center}

$^{*}$Alternation period is \texttt{ABORT\_LATCH\_COLOR\_INTERVAL\_MS} (see
§10 Parameters \& Flags).

\texttt{link\_connected} is set true on every received \cmd{HEARTBEAT} and
cleared after 15~s of silence (loss-of-signal watchdog).
\texttt{tank\_pressurized} is updated by \sw{Check\_Pressure\_Recording()}
using the live tank pressure reading; the threshold is
$>$\,\texttt{PRESSURE\_WARNING\_THRESHOLD} (15~PSI).

\subsection{Buzzer Behaviour}

\begin{center}
\begin{tabularx}{\textwidth}{@{}lY@{}}
\toprule
Condition & Buzzer Action \\
\midrule
Startup self-test            & 5 beeps (75~ms on, 75~ms gap) --- blocking, in \sw{setup()} \\
Button press confirmation    & 1 short beep via \sw{Beep()} (Button~A press) \\
Pressure alarm                & \textbf{Continuous tone} while \texttt{tank\_pressurized} = true \\
Engine start signal            & \textbf{Continuous tone} during the pre-flow phase of an RCU/LCO-sourced start (silenced automatically once propellant flow begins or the sequence ends) \\
Abort in progress              & \textbf{Continuous tone} while \texttt{emu\_state} = \texttt{ABORTING}; silenced once vented and latched \\
Fill beep                      & 100~ms timed beep every \texttt{FILL\_BEEP\_INTERVAL\_MS} while a fill session is active (filling or paused) \\
``I'm alive'' heartbeat        & 100~ms timed beep every \texttt{ALIVE\_BEEP\_INTERVAL\_MS}; suppressed during the pressure alarm, filling/paused, the engine-start signal or silent phases, \texttt{ENGINE\_RUNNING}, \texttt{ABORTING}, and \texttt{ABORT\_LATCHED} \\
\bottomrule
\end{tabularx}
\end{center}

The three continuous-tone conditions (pressure alarm, start signal, abort)
each track their own \texttt{\_active} flag so that when one turns off the
buzzer isn't stomped silent if another continuous-tone condition is still
active on the same tick. \sw{Execute\_Abort()} silences the buzzer alarm
immediately as part of its instant actions, independent of the sequenced
abort/vent procedure that follows.

\subsection{NeoPixel\_Indicator API}

\begin{center}
\begin{tabularx}{\textwidth}{@{}lY@{}}
\toprule
Method & Action \\
\midrule
\sw{begin()}                             & Initialise hardware; clear all pixels \\
\sw{on(color)}                           & Set all pixels solid; cancel any flash or timed mode \\
\sw{onRed()} / \sw{onGreen()} / \sw{onYellow()} & Convenience wrappers \\
\sw{flash(color, on\_ms, off\_ms)}       & Periodic flash indefinitely; cancelled by \sw{on()} or \sw{off()} \\
\sw{off()} / \sw{clear()}               & All pixels off; cancel flash/timed modes \\
\sw{update()}                            & Service timed shutoff and flash state machine --- call from loop() \\
\bottomrule
\end{tabularx}
\end{center}

%═══════════════════════════════════════════════════════════════════════════════
\section{Sensor Subsystems}
%═══════════════════════════════════════════════════════════════════════════════

\begin{center}
\begin{tabularx}{\textwidth}{@{}p{3.4cm}lp{2.4cm}Y@{}}
\toprule
Sensor & Class & ADC / Bus & Measurement \\
\midrule
Chamber pressure transducer& \sw{Pressure\_Transducer}   & ADS1015 ch1 & Combustion chamber pressure, 0--200~PSI \\
Tank pressure transducer   & \sw{Pressure\_Transducer}   & ADS1015 ch2 & N\textsubscript{2}O tank pressure, 0--1000~PSI \\
Battery voltage            & ---                         & ADS1015 ch3 & LiPo pack voltage (via resistor divider) \\
Airstart/LCO ignition sense & ---                        & ADS1015 ch0 & Ignition-enable signal (see §10 Parameters \& Flags, \texttt{LCO\_SENSE\_*}) \\
Thermocouple 1 (TC1)       & \sw{MAX31856}               & SPI         & Read by \cmd{READ\_TEMP1} (location TBD) -- all temperature sensing on this board; there is no separate onboard-ADC temperature path \\
Thermocouple 2 (TC2)       & \sw{MAX31856}               & SPI         & Read by \cmd{READ\_TEMP2} (location TBD) \\
\bottomrule
\end{tabularx}
\end{center}

\subsection{Pressure Transducer Signal Chain}

Each pressure transducer produces a ratiometric voltage output over
0.5--4.5~V corresponding to 0--full-scale PSI.  A voltage divider on the
A/D daughter board scales this to the ADS1015 input range:

\[
V_\text{ADC} = V_\text{transducer} \times \frac{3.3}{4.5}
\]

\begin{center}
\begin{tabular}{@{}lrrl@{}}
\toprule
Condition & Transducer (V) & After divider (V) & ADS1015 count (GAIN\_ONE) \\
\midrule
0 PSI (zero offset)  & 0.500 & 0.367 & $\approx$ \textbf{184} \\
Full-scale PSI       & 4.500 & 3.300 & $\approx$ \textbf{1657} \\
\bottomrule
\end{tabular}
\end{center}

ADS1015 at \texttt{GAIN\_ONE} ($\pm$4.096~V): $\approx$0.001991~V/count.
The two-point linear \sw{map()} in \sw{Pressure\_Sensor\_Average()} implicitly
performs all three conversion steps (undo divider, subtract 0.5~V offset,
scale 0--4~V to PSI) without explicit voltage arithmetic.

\begin{center}
\begin{tabular}{@{}lll@{}}
\toprule
Constant & Value & Meaning \\
\midrule
\sw{PRESSURE\_MIN\_READING} & 184  & ADS1015 counts at 0~PSI (0.500~V transducer) \\
\sw{PRESSURE\_MAX\_READING} & 1657 & ADS1015 counts at full-scale PSI (4.500~V transducer) \\
\sw{CHAMBER\_MAX\_PSI}      & 200  & Chamber transducer full-scale (PSI) \\
\sw{TANK\_MAX\_PSI}         & 1000 & Tank transducer full-scale (PSI) \\
\bottomrule
\end{tabular}
\end{center}

\subsection{Continuous Pressure Recording}

Continuous recording is started automatically by \cmd{BEGIN\_FILL} and may
also be started manually with \cmd{START\_PRESSURE\_RECORDING}. Every
\texttt{PRESSURE\_RECORDING\_INTERVAL\_MS} (250~ms) -- unless a burn
recording (below) currently owns the data file -- the EMU logs one row
with: elapsed time, battery voltage, tank and chamber pressure (PSI),
both thermocouple readings (\textdegree F), Ox/Fuel valve position (\%),
QR separation state, and igniter on/off.

\begin{sloppypar}
Recording auto-stops when tank pressure remains at or below
\texttt{PRESSURE\_WARNING\_THRESHOLD} (15~PSI) for
\texttt{PRESSURE\_AMBIENT\_COUNT} (4) consecutive readings, or when
\cmd{STOP\_PRESSURE\_RECORDING} is received. \cmd{RESET} does \emph{not}
stop it -- recording deliberately continues through an abort (see §3.4)
and stops on its own once pressure settles.
\end{sloppypar}

\subsection{Burn Data Recording}

A separate, higher-rate capture mechanism records propellant flow/burn
events to their own \texttt{PT\_N.txt} file, independent of the continuous
pressure log above. \sw{Arm\_Engine\_Start()} arms it automatically at
the start of every engine start sequence; \sw{Execute\_Abort()} stops and
writes out whatever was captured so far.

\begin{sloppypar}
While armed, \sw{Check\_Burn\_Recording()} (every \sw{loop()}) samples
both tared pressures and both thermocouples at \texttt{BURN\_SAMPLE\_INTERVAL\_MS}
(100~ms, 10~Hz) into a \sw{Burn\_Buffer} ring buffer holding
\texttt{BURN\_PRE\_TRIGGER\_MS} of history before the trigger and up to
\texttt{BURN\_POST\_TRIGGER\_DURATION\_MS} after it. The trigger fires once
chamber pressure has stayed \texttt{BURN\_TRIGGER\_DELTA\_PSI} above tared
ambient for \texttt{BURN\_TRIGGER\_SUSTAIN\_MS}; if it never fires, a
\texttt{BURN\_NEVER\_TRIGGERED\_TIMEOUT\_MS} safety net completes the
recording anyway (see §10 Parameters \& Flags for all five values). Once
complete, the full pre/post-trigger buffer is written to SD in one burst
and the file is closed.
\end{sloppypar}

%═══════════════════════════════════════════════════════════════════════════════
\section{SD Logging}
%═══════════════════════════════════════════════════════════════════════════════

Log files are stored in the root of the SD card as \texttt{Log\_N.txt}
(N is auto-incremented to avoid overwriting existing files).
Each record is time-stamped (PCF8523) and records command receipt, valve
movements, sensor readings, and fault events.

%═══════════════════════════════════════════════════════════════════════════════
\section{Parameters \& Flags}
%═══════════════════════════════════════════════════════════════════════════════

This section collects the tunable operational parameters and feature flags
compiled into this firmware image, together with the source file where each
is defined. It is meant as a single reference point when tracing or changing
a timing, threshold, or calibration value. Pin/channel wiring assignments,
LoRa/CAN command opcodes, and packet byte layouts are documented separately
(the Hardware Overview pin map and the Command Dispatch/Command Table, §5, above) and are
not repeated here. Parameters marked \emph{(shared)} are defined in
\sw{KJO\_\allowbreak Shared\_\allowbreak Libraries} and must agree with the same constant used on
the other board(s) that consume it.

\subsubsection*{Valve Servo Timing \& Calibration}

\begingroup
\footnotesize
\begin{longtable}{@{}p{3.2cm}p{2.4cm}p{6.2cm}p{2.6cm}@{}}
\toprule
Parameter & Value & Description & Source File \\
\midrule
\texttt{SERVO\_\allowbreak FREQUENCY} & 1100~Hz & Non-standard PWM frequency (vs.\ 50~Hz) chosen to compress the usable PWM count range and improve position resolution over the servo's $\sim$1.1-turn travel. & \texttt{KJO\_\allowbreak Valve.h} \\
\texttt{SERVO\_\allowbreak RANGE\_\allowbreak MIN} & 500~$\mu$s & Pulse width at the servo's lower mechanical travel limit. & \texttt{KJO\_\allowbreak Valve.h} \\
\texttt{SERVO\_\allowbreak RANGE\_\allowbreak MAX} & 2500~$\mu$s & Pulse width at the servo's upper mechanical travel limit. & \texttt{KJO\_\allowbreak Valve.h} \\
\texttt{SERVO\_\allowbreak FULL\_\allowbreak RANGE\_\allowbreak DEGREES} & 1800\textdegree & Total mechanical travel of the 5-turn servo used for position-mapping math. & \texttt{KJO\_\allowbreak Valve.h} \\
\texttt{SERVO\_\allowbreak PINION\_\allowbreak TEETH} & 16 & Tooth count of the servo output pinion gear, used in valve gearing geometry. & \texttt{KJO\_\allowbreak Valve.h} \\
\texttt{VALVE\_\allowbreak GEAR\_\allowbreak TEETH} & 60 & Tooth count of the valve actuator gear; combined with pinion teeth sets the 60:16 gear ratio. & \texttt{KJO\_\allowbreak Valve.h} \\
\texttt{VALVE\_\allowbreak TURN\_\allowbreak DEGREES} & 90\textdegree & Valve rotational travel from full open to full closed. & \texttt{KJO\_\allowbreak Valve.h} \\
\texttt{OX\_\allowbreak VALVE\_\allowbreak PWM\_\allowbreak OPEN} & 2200 & 12-bit PWM count for the Ox valve's fully-open servo position. & \texttt{KJO\_\allowbreak Valve.h} \\
\texttt{OX\_\allowbreak VALVE\_\allowbreak PWM\_\allowbreak CLOSE} & 3800 & 12-bit PWM count for the Ox valve's fully-closed servo position. & \texttt{KJO\_\allowbreak Valve.h} \\
\texttt{OX\_\allowbreak VALVE\_\allowbreak POS\_\allowbreak OPEN} & 1141 & ADC count from the Ox valve position pot when fully open. & \texttt{KJO\_\allowbreak Valve.h} \\
\texttt{OX\_\allowbreak VALVE\_\allowbreak POS\_\allowbreak CLOSED} & 2898 & ADC count from the Ox valve position pot when fully closed. & \texttt{KJO\_\allowbreak Valve.h} \\
\texttt{OX\_\allowbreak VALVE\_\allowbreak BALL\_\allowbreak DIAMETER} & 0.611~in & Ball outer diameter used in the Ox valve's angle-based flow-position geometry calculation. & \texttt{KJO\_\allowbreak Valve.h} \\
\texttt{OX\_\allowbreak VALVE\_\allowbreak BORE\_\allowbreak DIAMETER} & 0.357~in & Bore diameter used in the Ox valve's flow-position geometry calculation. & \texttt{KJO\_\allowbreak Valve.h} \\
\texttt{OX\_\allowbreak VALVE\_\allowbreak THROAT\_\allowbreak DIAMETER} & 0.373~in & Throat inner diameter used in the Ox valve's flow-position geometry calculation. & \texttt{KJO\_\allowbreak Valve.h} \\
\texttt{FUEL\_\allowbreak VALVE\_\allowbreak PWM\_\allowbreak OPEN} & 2200 & 12-bit PWM count for the Fuel valve's fully-open servo position. & \texttt{KJO\_\allowbreak Valve.h} \\
\texttt{FUEL\_\allowbreak VALVE\_\allowbreak PWM\_\allowbreak CLOSE} & 3800 & 12-bit PWM count for the Fuel valve's fully-closed servo position. & \texttt{KJO\_\allowbreak Valve.h} \\
\texttt{FUEL\_\allowbreak VALVE\_\allowbreak POS\_\allowbreak OPEN} & 1108 & ADC count from the Fuel valve position pot when fully open. & \texttt{KJO\_\allowbreak Valve.h} \\
\texttt{FUEL\_\allowbreak VALVE\_\allowbreak POS\_\allowbreak CLOSED} & 2786 & ADC count from the Fuel valve position pot when fully closed. & \texttt{KJO\_\allowbreak Valve.h} \\
\texttt{FUEL\_\allowbreak VALVE\_\allowbreak BALL\_\allowbreak DIAMETER} & 0.611~in & Ball outer diameter used in the Fuel valve's flow-position geometry calculation. & \texttt{KJO\_\allowbreak Valve.h} \\
\texttt{FUEL\_\allowbreak VALVE\_\allowbreak BORE\_\allowbreak DIAMETER} & 0.357~in & Bore diameter used in the Fuel valve's flow-position geometry calculation. & \texttt{KJO\_\allowbreak Valve.h} \\
\texttt{FUEL\_\allowbreak VALVE\_\allowbreak THROAT\_\allowbreak DIAMETER} & 0.373~in & Throat inner diameter used in the Fuel valve's flow-position geometry calculation. & \texttt{KJO\_\allowbreak Valve.h} \\
\texttt{VALVE\_\allowbreak MOVE\_\allowbreak TIMEOUT\_\allowbreak MS} (alias \texttt{MOVE\_\allowbreak TIMEOUT}) & 2500~ms & Maximum time allowed for a commanded valve move to complete before the watchdog gives up; shared with GSEMU since both use the identical servo/gear actuator. & \texttt{KJO\_\allowbreak Valve\_\allowbreak Timing.h} (shared) \\
\texttt{VALVE\_\allowbreak POSITION\_\allowbreak AVERAGING\_\allowbreak COUNT} & 4 & Number of ADC reads averaged for each valve position sample. & \texttt{KJO\_\allowbreak Valve.h} \\
\texttt{VALVE\_\allowbreak MOVING\_\allowbreak THRESHOLD} & 2~counts & Minimum encoder delta between samples required to consider the valve still moving. & \texttt{KJO\_\allowbreak Valve.h} \\
\texttt{VALVE\_\allowbreak POSITION\_\allowbreak DEAD\_\allowbreak BAND} & 15~counts & Position tolerance used to decide whether a valve has reached its open/closed target. & \texttt{KJO\_\allowbreak Valve.h} \\
\texttt{MOVE\_\allowbreak DETECT\_\allowbreak WAIT} & 30~ms & Minimum interval between successive encoder samples inside \sw{isMoving()}. & \texttt{KJO\_\allowbreak Valve.h} \\
\texttt{MOVE\_\allowbreak AFTER\_\allowbreak START\_\allowbreak WAIT} & 60~ms & Blocking delay after issuing a servo PWM command, giving the servo time to begin responding before motion is sampled. & \texttt{KJO\_\allowbreak Valve.h} \\
\texttt{USE\_\allowbreak MAPPING} & \texttt{true} & Selects angle-geometry-based valve position mapping over simple linear PWM interpolation. & \texttt{KJO\_\allowbreak System\_\allowbreak Config.h} (shared) \\
\end{longtable}
\endgroup

\subsubsection*{Pressure Sensing \& Calibration}

\begingroup
\footnotesize
\begin{longtable}{@{}p{3.2cm}p{2.4cm}p{6.2cm}p{2.6cm}@{}}
\toprule
Parameter & Value & Description & Source File \\
\midrule
\texttt{PRESSURE\_\allowbreak AVERAGING\_\allowbreak COUNT} & 4 & Number of ADS1015 reads averaged per pressure measurement. & \texttt{KJO\_\allowbreak Pressure.h} \\
\texttt{CHAMBER\_\allowbreak PRESSURE\_\allowbreak SLOPE} & 0.137533 & Bench-calibrated linear-fit slope (PSI per ADC count) for the chamber pressure transducer. & \texttt{KJO\_\allowbreak Pressure.h} \\
\texttt{CHAMBER\_\allowbreak PRESSURE\_\allowbreak INTERCEPT} & $-$25.012515 & Bench-calibrated linear-fit intercept (PSI at 0 counts) for the chamber pressure transducer. & \texttt{KJO\_\allowbreak Pressure.h} \\
\texttt{TANK\_\allowbreak PRESSURE\_\allowbreak SLOPE} & 0.688427 & Bench-calibrated linear-fit slope (PSI per ADC count) for the Ox tank pressure transducer. & \texttt{KJO\_\allowbreak Pressure.h} \\
\texttt{TANK\_\allowbreak PRESSURE\_\allowbreak INTERCEPT} & $-$124.131326 & Bench-calibrated linear-fit intercept (PSI at 0 counts) for the Ox tank pressure transducer. & \texttt{KJO\_\allowbreak Pressure.h} \\
\texttt{TANK\_\allowbreak MAX\_\allowbreak PSI} & 1000~PSI & Full-scale rating of the tank pressure transducer, used for header/documentation purposes. & \texttt{KJO\_\allowbreak Pressure.h} \\
\texttt{CHAMBER\_\allowbreak MAX\_\allowbreak PSI} & 200~PSI & Full-scale rating of the chamber pressure transducer, used for header/documentation purposes. & \texttt{KJO\_\allowbreak Pressure.h} \\
\texttt{PRESSURE\_\allowbreak WARNING\_\allowbreak THRESHOLD} & 15~PSI & Tank pressure above which the red indicator and pressure-alarm buzzer activate, and below which recording auto-stops. & \texttt{KJO\_\allowbreak Pressure.h} \\
\texttt{PRESSURE\_\allowbreak CHECK\_\allowbreak INTERVAL\_\allowbreak MS} & 250~ms & Poll interval for the unconditional tank-pressure safety check (\sw{Check\_\allowbreak Pressure\_\allowbreak Warning()}). & \texttt{KJO\_\allowbreak Pressure.h} \\
\texttt{ABORT\_\allowbreak VENTED\_\allowbreak THRESHOLD\_\allowbreak PSI} & 5~PSI & Tank gauge pressure below which an in-progress Abort's vent-wait step considers the tank vented and latches. & \texttt{KJO\_\allowbreak Pressure.h} \\
\texttt{MAX\_\allowbreak OX\_\allowbreak TANK\_\allowbreak PRESSURE} & 950~PSI & Sustained Ox-tank gauge pressure threshold that triggers an automatic Abort. & \texttt{KJO\_\allowbreak Pressure.h} \\
\texttt{OX\_\allowbreak TANK\_\allowbreak PRESSURE\_\allowbreak TIME\_\allowbreak MS} & 1000~ms & How long the Ox tank must remain at/above \texttt{MAX\_\allowbreak OX\_\allowbreak TANK\_\allowbreak PRESSURE} before the automatic abort fires. & \texttt{KJO\_\allowbreak Pressure.h} \\
\texttt{AD\_\allowbreak BASE\_\allowbreak SCALE} & 0.00199144777 V/count & ADS1015 volts-per-count scale factor at PGA gain 1, used for battery-voltage conversion. & \texttt{KJO\_\allowbreak Analog.h} \\
\texttt{LIPO\_\allowbreak SCALE} & 2.79111 & Resistor-divider scale factor converting the ADS1015 LiPo channel reading to actual battery voltage. & \texttt{KJO\_\allowbreak Analog.h} \\
\texttt{AIRSTART\_\allowbreak THRESHOLD} (alias of \texttt{LCO\_\allowbreak SENSE\_\allowbreak THRESHOLD}) & 251~counts & ADC-count threshold above which the Airstart/LCO sense channel is considered triggered ($\sim$0.5~V post-divider). & \texttt{KJO\_\allowbreak LCO\_\allowbreak Sense.h} (shared) \\
\end{longtable}
\endgroup

\subsubsection*{Thermocouple Sensing}

\begingroup
\footnotesize
\begin{longtable}{@{}p{3.2cm}p{2.4cm}p{6.2cm}p{2.6cm}@{}}
\toprule
Parameter & Value & Description & Source File \\
\midrule
\texttt{TEMPERATURE\_\allowbreak TC\_\allowbreak AVERAGING\_\allowbreak COUNT} & 4 & Number of thermocouple readings averaged per measurement (declared but not yet wired into a read path). & \texttt{KJO\_\allowbreak Thermocouple.h} \\
\end{longtable}
\endgroup

\subsubsection*{Engine Start Sequence Timing}

\begingroup
\footnotesize
\begin{longtable}{@{}p{3.2cm}p{2.4cm}p{6.2cm}p{2.6cm}@{}}
\toprule
Parameter & Value & Description & Source File \\
\midrule
\texttt{IGNITER\_\allowbreak FIRE\_\allowbreak DURATION\_\allowbreak MS} & 2000~ms & Duration the igniter relay stays energized to reliably light the solid starter motor. & \texttt{main.cpp} \\
\texttt{START\_\allowbreak DWELL\_\allowbreak MS} & 2000~ms & Fixed pause after confirmed umbilical separation before propellant flow begins. & \texttt{main.cpp} \\
\texttt{FUEL\_\allowbreak FLOW\_\allowbreak DELAY\_\allowbreak MS} & 500~ms & Placeholder transit-time allowance after Fuel valve confirms open, before proceeding (pending cold-flow calibration). & \texttt{main.cpp} \\
\texttt{OX\_\allowbreak FLOW\_\allowbreak DELAY\_\allowbreak MS} & 200~ms & Placeholder transit-time allowance after Ox valve confirms open, before igniter fires (pending cold-flow calibration). & \texttt{main.cpp} \\
\texttt{QR\_\allowbreak RELEASE\_\allowbreak TIMEOUT\_\allowbreak MS} & 3000~ms & Maximum time to wait for physical umbilical separation after commanding QR release, before aborting the start sequence. & \texttt{main.cpp} \\
\texttt{ABORT\_\allowbreak VALVE\_\allowbreak CONFIRM\_\allowbreak TIMEOUT\_\allowbreak MS} & \texttt{MOVE\_\allowbreak TIMEOUT} + 100~ms (2600~ms) & Watchdog for each Abort-sequence valve-close/open confirmation step. & \texttt{main.cpp} \\
\texttt{EMU\_\allowbreak GSEMU\_\allowbreak TIMEOUT\_\allowbreak MS} & 200~ms & CAN round-trip timeout for ordinary EMU-to-GSEMU relay requests (no chaining through GSEMU). & \texttt{main.cpp} \\
\texttt{EMU\_\allowbreak LCMU\_\allowbreak TIMEOUT\_\allowbreak MS} & 200~ms & CAN round-trip timeout for ordinary EMU-to-LCMU relay requests. & \texttt{main.cpp} \\
\texttt{EMU\_\allowbreak GSEMU\_\allowbreak FILL\_\allowbreak VALVE\_\allowbreak MOVE\_\allowbreak TIMEOUT\_\allowbreak MS} & \texttt{MOVE\_\allowbreak TIMEOUT} + 300~ms (2800~ms) & CAN timeout for GSEMU Fill-valve open/close requests, sized to cover GSEMU's own valve-move watchdog plus round-trip margin. & \texttt{main.cpp} \\
\texttt{Motor\_\allowbreak Ox\_\allowbreak Fill\_\allowbreak Target\_\allowbreak Lbm()} values (L/M/N/O) & 7.026 / 14.053 / 28.106 / 58.219~lbm & Per-motor-designation Ox fill target weight EMU's fill session compares polled LCMU readings against. & \texttt{main.cpp} \\
\texttt{SELECTED\_\allowbreak MOTOR} & \texttt{Motor\_\allowbreak Designation::M} & Selects which motor's Ox fill target is active for the current test article; changed per test. & \texttt{main.cpp} \\
\end{longtable}
\endgroup

\subsubsection*{Abort / Indicator / Buzzer Timing}

\begingroup
\footnotesize
\begin{longtable}{@{}p{3.2cm}p{2.4cm}p{6.2cm}p{2.6cm}@{}}
\toprule
Parameter & Value & Description & Source File \\
\midrule
\texttt{LINK\_\allowbreak LOS\_\allowbreak INTERVAL\_\allowbreak MS} & 15000~ms & Time without a received radio packet before the RCU link is declared lost. & \texttt{main.cpp} \\
\texttt{ALIVE\_\allowbreak BEEP\_\allowbreak INTERVAL\_\allowbreak MS} & 10000~ms & Period of the idle ``I'm alive'' buzzer beep. & \texttt{main.cpp} \\
\texttt{FILL\_\allowbreak BEEP\_\allowbreak INTERVAL\_\allowbreak MS} & 2000~ms & Period of the buzzer beep while a fill session is active (filling or paused), replacing the alive beep. & \texttt{main.cpp} \\
\texttt{ABORT\_\allowbreak LATCH\_\allowbreak COLOR\_\allowbreak INTERVAL\_\allowbreak MS} & 1000~ms & Red/yellow alternation period (per color) for the indicator once ABORT\_\allowbreak LATCHED. & \texttt{main.cpp} \\
\texttt{BOOT\_\allowbreak RETRY\_\allowbreak BACKOFF\_\allowbreak MS} & 3000~ms & Delay before retrying a timed-out boot-time CAN presence ping (GSEMU/LCMU). & \texttt{main.cpp} \\
\texttt{PRESENCE\_\allowbreak PING\_\allowbreak TIMEOUT\_\allowbreak MS} & 200~ms & CAN round-trip timeout for the runtime GSEMU/LCMU presence ping. & \texttt{main.cpp} \\
\texttt{PRESENCE\_\allowbreak PING\_\allowbreak INTERVAL\_\allowbreak MS} & 2000~ms & Interval between successive full GSEMU+LCMU presence-ping cycles. & \texttt{main.cpp} \\
\texttt{FILL\_\allowbreak POLL\_\allowbreak INTERVAL\_\allowbreak MS} & 200~ms & Poll interval for the background LCMU weight query during an active fill. & \texttt{main.cpp} \\
\end{longtable}
\endgroup

\subsubsection*{Continuous Sensor Recording / Burn Data Capture}

\begingroup
\footnotesize
\begin{longtable}{@{}p{3.2cm}p{2.4cm}p{6.2cm}p{2.6cm}@{}}
\toprule
Parameter & Value & Description & Source File \\
\midrule
\texttt{PRESSURE\_\allowbreak RECORDING\_\allowbreak INTERVAL\_\allowbreak MS} & 250~ms & Sample period for continuous pressure/temperature recording while active. & \texttt{main.cpp} \\
\texttt{PRESSURE\_\allowbreak AMBIENT\_\allowbreak COUNT} & 4 & Consecutive below-threshold pressure samples required to confirm ambient and auto-stop recording. & \texttt{main.cpp} \\
\texttt{BURN\_\allowbreak TRIGGER\_\allowbreak DELTA\_\allowbreak PSI} & 5.0~PSI & Chamber gauge pressure rise above tared ambient that starts the burn-trigger sustain timer. & \texttt{main.cpp} \\
\texttt{BURN\_\allowbreak TRIGGER\_\allowbreak SUSTAIN\_\allowbreak MS} & 1000~ms & How long chamber pressure must stay above \texttt{BURN\_\allowbreak TRIGGER\_\allowbreak DELTA\_\allowbreak PSI} before the burn recording is considered triggered. & \texttt{main.cpp} \\
\texttt{BURN\_\allowbreak SAMPLE\_\allowbreak INTERVAL\_\allowbreak MS} & 100~ms (10~Hz) & Sample rate for burn-data capture. & \texttt{main.cpp} \\
\texttt{BURN\_\allowbreak PRE\_\allowbreak TRIGGER\_\allowbreak MS} & 10000~ms & Amount of pre-trigger history retained in the burn ring buffer. & \texttt{main.cpp} \\
\texttt{BURN\_\allowbreak POST\_\allowbreak TRIGGER\_\allowbreak DURATION\_\allowbreak MS} & 30000~ms & Fixed capture duration after the burn trigger fires. & \texttt{main.cpp} \\
\texttt{BURN\_\allowbreak NEVER\_\allowbreak TRIGGERED\_\allowbreak TIMEOUT\_\allowbreak MS} & 60000~ms & Safety-net duration after arming; if the burn never triggers, recording completes anyway. & \texttt{main.cpp} \\
\texttt{BURN\_\allowbreak BUFFER\_\allowbreak MAX\_\allowbreak POST\_\allowbreak SAMPLES} & 300 (30~s @ 10~Hz) & Fixed capacity of the post-trigger burn sample buffer. & \texttt{KJO\_\allowbreak Burn\_\allowbreak Buffer.h} \\
\texttt{BURN\_\allowbreak BUFFER\_\allowbreak MAX\_\allowbreak PRE\_\allowbreak SAMPLES} & 100 (10~s @ 10~Hz) & Fixed capacity of the pre-trigger burn ring buffer. & \texttt{KJO\_\allowbreak Burn\_\allowbreak Buffer.h} \\
\texttt{BURN\_\allowbreak BUFFER\_\allowbreak NUM\_\allowbreak CHANNELS} & 4 & Number of sensor channels stored per burn sample (tank PSI, chamber PSI, Temp1, Temp2). & \texttt{KJO\_\allowbreak Burn\_\allowbreak Buffer.h} \\
\end{longtable}
\endgroup

\subsubsection*{ADC Diagnostic Mode (bench-only, compiled out by default)}

\begingroup
\footnotesize
\begin{longtable}{@{}p{3.2cm}p{2.4cm}p{6.2cm}p{2.6cm}@{}}
\toprule
Parameter & Value & Description & Source File \\
\midrule
\texttt{ADS\_\allowbreak DIAG\_\allowbreak LSB\_\allowbreak V} & 4.096/2048 V/count & ADS1015 GAIN\_\allowbreak ONE LSB voltage used to back-calculate connector voltage during bench diagnostics. & \texttt{main.cpp} \\
\texttt{ADS\_\allowbreak DIAG\_\allowbreak DIV\_\allowbreak INV} & 4.5/3.3 & Voltage-divider inverse ratio for the chamber channel diagnostic readback. & \texttt{main.cpp} \\
\texttt{ADS\_\allowbreak DIAG\_\allowbreak DIV\_\allowbreak INV\_\allowbreak TANK} & $(4.5\times3.290)/(3.3\times3.267) \approx 1.37324$ & Bench-trimmed voltage-divider inverse ratio for the tank channel diagnostic readback. & \texttt{main.cpp} \\
\texttt{ADC\_\allowbreak DIAG\_\allowbreak INTERVAL\_\allowbreak MS} & 2000~ms & How often the ADC diagnostic reading is printed to the OLED. & \texttt{main.cpp} \\
\end{longtable}
\endgroup

\subsubsection*{Radio Communication}

\begingroup
\footnotesize
\begin{longtable}{@{}p{3.2cm}p{2.4cm}p{6.2cm}p{2.6cm}@{}}
\toprule
Parameter & Value & Description & Source File \\
\midrule
\texttt{RCS\_\allowbreak CENTER\_\allowbreak FREQUENCY} & 917.117~MHz & LoRa center frequency; must match between EMU and RCU radios to communicate. & \texttt{KJO\_\allowbreak Radio.h} (shared) \\
\texttt{TRANSMIT\_\allowbreak POWER} & 5~dBm & LoRa transmit power setting (valid range 5--23~dBm via PA\_\allowbreak BOOST). & \texttt{KJO\_\allowbreak Radio.h} (shared) \\
\texttt{LINK\_\allowbreak MARGINAL\_\allowbreak THRESHOLD} & $-$90~dBm & RSSI below which the link-quality indicator shows MARGINAL rather than OK. & \texttt{KJO\_\allowbreak Radio.h} (shared) \\
\texttt{MAX\_\allowbreak PAYLOAD\_\allowbreak LENGTH} & 64~bytes & Maximum LoRa packet payload size EMU/RCU buffers allocate (headroom under RadioHead's 251-byte ceiling). & \texttt{KJO\_\allowbreak Command\_\allowbreak Defs.h} (shared) \\
\texttt{MAX\_\allowbreak COMMAND\_\allowbreak CODE} & 39 & Highest valid LoRa command code, used to range-check incoming packets. & \texttt{KJO\_\allowbreak Command\_\allowbreak Defs.h} (shared) \\
\end{longtable}
\endgroup

\subsubsection*{CAN Bus Communication Timing}

\begingroup
\footnotesize
\begin{longtable}{@{}p{3.2cm}p{2.4cm}p{6.2cm}p{2.6cm}@{}}
\toprule
Parameter & Value & Description & Source File \\
\midrule
\texttt{CAN\_\allowbreak BAUDRATE} & 1{,}000{,}000~bps & CAN bus bit rate shared by EMU/GSEMU/LCMU; raised from 250~kbps, safely within spec for the installation's short (16~ft max) bus runs. & \texttt{KJO\_\allowbreak CAN\_\allowbreak Command\_\allowbreak Defs.h} (shared) \\
\texttt{CAN\_\allowbreak MAX\_\allowbreak COMMAND\_\allowbreak CODE} & 23 & Highest valid CAN command code, used to bounds-check inbound frames before name-table lookup. & \texttt{KJO\_\allowbreak CAN\_\allowbreak Command\_\allowbreak Defs.h} (shared) \\
\end{longtable}
\endgroup

\subsubsection*{SD Logging}

\begingroup
\footnotesize
\begin{longtable}{@{}p{3.2cm}p{2.4cm}p{6.2cm}p{2.6cm}@{}}
\toprule
Parameter & Value & Description & Source File \\
\midrule
\texttt{LOG\_\allowbreak FILE\_\allowbreak NAME\_\allowbreak BASE} & \texttt{"Log\_\allowbreak "} & Base filename prefix for sequentially-numbered SD log files (e.g.\ \texttt{Log\_\allowbreak 0.txt}). & \texttt{KJO\_\allowbreak Logging.h} \\
\end{longtable}
\endgroup

%═══════════════════════════════════════════════════════════════════════════════
\section{Version History}
%═══════════════════════════════════════════════════════════════════════════════
\begin{longtable}{@{}p{1.4cm}p{2.2cm}p{11.5cm}@{}}
\toprule
Version & Date & Notes \\
\midrule
0.1 & March 2026 & Initial release. \\
0.2 & March 2026 & Add umbilical QR release subsystem (QR\_Master, QR\_RELEASE command, AUX\_IO\_3 protocol). \\
0.3 & March 2026 & Level-based QR protocol (release=LOW, hold=HIGH); add RELEASE\_STATE\_EIO,
                   LAUNCH\_ENABLE\_EIO, REMOTE\_START\_EIO; Button~A reassigned to hold-to-open
                   Dump valve; QR\_Master API updated. \\
0.4 & March 2026 & Full command dispatch table; ENGINE\_STARTING state added to state machine;
                   7-step engine start sequence documented; new §7 Buzzer and Visual Indicator:
                   5-beep startup, NeoPixel state machine (5 states/priorities), continuous
                   pressure-alarm buzzer, 10~s alive beep, alarm silence on start/abort. \\
0.5 & March 2026 & Valve servo calibration: added §3 servo geometry subsection (gear ratio
                   3.75:1, 1100~Hz PCA9685, $\approx$1690 theoretical counts for 90\textdegree)
                   and empirical per-valve calibration table (all four valves re-calibrated on
                   bench). Pressure transducer fix: corrected ADS1015 channel assignments
                   (chamber = ch1/200~PSI, tank = ch2/1000~PSI), corrected sensor table
                   (temperature sensors use onboard ADC A0/A5, not ADS1015), added §7 signal
                   chain subsection documenting voltage divider and calibration endpoint counts
                   (184/1657). \\
0.6 & March 2026 & SD log files moved to SD card root (\texttt{Log\_N.txt}); subdirectory path
                   constants (\texttt{LOG\_FILE\_FOLDER}, \texttt{LOG\_FILE\_DIR}) removed from
                   \sw{KJO\_Logging.h} and \sw{KJO\_EMU\_Display}. \\
0.7 & July 2026  & Dump valve physically removed from the vehicle; pressurant is now dumped
                   through the combustion chamber via the Ox valve. \cmd{DUMP}/\cmd{STOP\_DUMP}
                   retained but now target the Ox valve; \sw{Execute\_Abort()} closes Fill/Fuel
                   and opens Ox (instead of closing all valves); Button~A reassigned to
                   hold-to-open the Ox valve; \texttt{dump\_pct} removed from
                   \texttt{Heartbeat\_Response\_Packet}. \\
0.8 & September 2026 & Add new §10 Parameters \& Flags reference section: tunable timing,
                   threshold, and calibration constants across the valve, pressure,
                   engine-start, indicator, burn-recording, radio, and CAN subsystems,
                   with source-file attribution for each. \\
0.9 & September 2026 & Corrected stale wired-GPIO descriptions left over from the CAN
                   redesign: removed \S3.5 Master\_Valve Class (the class no longer
                   exists); rewrote \S4 as ``Fill Valve CAN Protocol'' describing the
                   actual \cmd{CAN\_OPEN\_FILL\_VALVE}/\cmd{CAN\_CLOSE\_FILL\_VALVE}
                   mechanism; rewrote \S6 Umbilical Quick-Release to describe
                   \sw{QR\_Master} as a sense-only reader and \cmd{CAN\_QR\_RELEASE} as
                   the real release path; corrected 5 of 6 \sig{AUX\_IO} pin-map rows
                   (only \sig{BUTTON\_A\_EIO}'s row and the umbilical/link-sense
                   functions survive; Fill valve CMD/STATUS and the old QR/launch/
                   remote-start lines are retired); added new \S3.5 Local Button
                   Functions documenting all three front-panel buttons (only Button~A
                   is still functional); corrected two stale Command Table rows
                   (\cmd{BEGIN\_FILL}, \cmd{QR\_RELEASE}). \\
1.0 & September 2026 & Comprehensive doc-vs-code audit: added the missing
                   \texttt{ABORTING}/\texttt{ABORT\_LATCHED} states and a new \S3.4 Abort
                   Sequence (5 states/7 steps were entirely undocumented); rewrote \S3.3 from
                   the stale 7-step, igniter-first ordering to the actual 10-step sequence;
                   rebuilt the Command Table from 22 to all 39 defined commands, fixing the
                   nonexistent \cmd{ENABLE\_LCO\_CONTROL} row (renamed \cmd{ARM\_LAUNCH}),
                   stale \cmd{READ\_TANK\_TEMP}/\cmd{READ\_CHAMBER\_TEMP} names, and the
                   \cmd{HEARTBEAT} row's fill-position claim; rebuilt \S2.2 Main Loop Control
                   Flow from 9 boxes (two -- \sw{Check\_Remote\_Start()},
                   \sw{Check\_Radio()} -- naming nonexistent functions) to the actual 20-step
                   \sw{loop()}; rebuilt \S7.2/\S7.3's indicator/buzzer tables from 5 to 10
                   states, fixing a wrong color and wrong flash timings; removed two fictional
                   \sw{Temperature\_Transducer} sensor rows and a nonexistent \sw{KJO\_Temperature}
                   module (all temperature sensing is via the two thermocouples); removed a
                   Fill-valve calibration row for a valve that isn't on this board; added the
                   MCP2515 CAN controller to the Hardware Overview (previously omitted
                   entirely); corrected the ADC channel~0 assignment; expanded the Module
                   Summary; documented the previously-unmentioned Burn Data Recording
                   subsystem (new \S8.3). \\
\bottomrule
\end{longtable}

\end{document}
